% Part: applied-modal-logics
% Chapter: epistemic-logic
% Section: public-announcement-logic-semantics

\documentclass[../../../include/open-logic-section]{subfiles}

\begin{document}

\olfileid{aml}{el}{psm}

\olsection{公共宣告逻辑的语义}

公共宣告逻辑的关系模型与认知逻辑中的关系模型完全一样。不过，公共宣告算子的语义则是全新的。

\begin{defn}\ollabel{defn:mmodels} 公式~$!A$ 在一个~$\mModel M = \tuple{W, R, V}$
  的世界~$w$ 处的\emph{真值}，记作：
  $\mSat{M}{!A}[w]$，归纳定义如下：
  \begin{enumerate}
  \tagitem{prvFalse}{\indcase{!A}{\lfalse}{永不 $\mSat{M}{\lfalse}[w]$}.}{}
  \tagitem{prvTrue}{\indcase{!A}{\ltrue}{总是 $\mSat{M}{\ltrue}[w]$}.}{}
  \item $\mSat{M}{p}[w]$ 当且仅当 $w \in V(p)$
  \tagitem{prvNot}{\indcase{!A}{\lnot !B}{$\mSat{M}{\indfrm}[w]$ 当且仅当
    $\mSat/{M}{!B}[w]$}.}{}
  \tagitem{prvAnd}{\indcase{!A}{(!B \land !C)}{$\mSat{M}{\indfrm}[w]$ 当且仅当
    $\mSat{M}{!B}[w]$ 且 $\mSat{M}{!C}[w]$}.}{}
  \tagitem{prvOr}{\indcase{!A}{(!B \lor !C)}{$\mSat{M}{\indfrm}[w]$ 当且仅当
    $\mSat{M}{!B}[w]$ 或 $\mSat{M}{!C}[w]$}（或两者皆成立）。}{}
  \tagitem{prvIf}{\indcase{!A}{(!B \lif !C)}{$\mSat{M}{\indfrm}[w]$ 当且仅当
    $\mSat/{M}{!B}[w]$ 或 $\mSat{M}{!C}[w]$}.}{}
  \tagitem{prvIff}{\indcase{!A}{(!B \liff !C)}{$\mSat{M}{\indfrm}[w]$ 当且仅当
    要么$\mSat{M}{!B}[w]$且$\mSat{M}{!C}[w]$，要么既不$\mSat{M}{!B}[w]$也不$\mSat{M}{!C}[w]$}.}{}
  \item\ollabel{defn:sub:mmodels-box}
    \indcase{!A}{\Knows_a !B}{对所有满足 $R_a ww'$ 的 $w' \in W$，$\mSat{M}{\indfrm}[w]$ 当且仅当
    $\mSat{M}{!B}[w']$}
  \item\ollabel{defn:sub:mmodels-pal}
    \indcase{!A}{[!B] !C}{$\mSat{M}{\indfrm}[w]$ 当且仅当
    如果$\mSat{M}{!B}[w]$那么$\mSat{M \mid !B}{!C}[w]$}


  其中 $\mModel M \mid !B = \tuple{W', R', V'}$ 定义如下：
  \begin{enumerate}
    \item $W' = \Setabs{ u \in W}{\mSat{M}{!B}[u]}$。因此，$\mModel M \mid !B$ 的世界就是 $\mModel M$ 中 $!B$ 成立的那些世界。
    \item $R'_a = R_a \cap (W' \times W')$。每个主体的可及关系被简单地限制在 $W'$ 中剩余的世界上。
    \item $V'(p) = \Setabs{ u \in W'}{u \in V(p) }$。类似地，世界上的命题赋值保持不变，这体现了信息事件不会改变命题变元真值的思想。
  \end{enumerate}

  \end{enumerate}
\end{defn}

那么，公共宣告逻辑的独特之处在于，一个公式在~$\mModel M$ 中是否为真，有时只能通过参考 $\mModel M$ 自身之外的另一个模型来决定。

还需注意，我们的语义将宣告算子视为一个 $\Box$~算子（必然算子），因此如果一个公式~$!A$ 在某个世界上不能被如实地宣告，那么 $[!A]B$ 将在此处平凡地成立，就像所有 $\Box$ 公式在端点上都成立一样。

\begin{figure}
  \begin{center}
    \begin{tikzpicture}[modal]
      \node[world] (w1) [label={below left:$p, \lnot q$}]{$w_1$};
      \node[world] (w2) [left=of w1, label={left:$\lnot p, \lnot q$}]{$w_2$};
      \node[world] (w3) [above=of w1, label={right:$p, q$}] {$w_3$};
      \draw[<->] (w1) to [swap] node {$b$} (w2);
      \draw[->,loop below] (w1) to node {$a, b$} (w1);
      \draw[<->] (w1) to [swap] node {$a$} (w3);
      \draw[->,loop above] (w2) to node {$a, b$} (w2) ;
      \draw[->,loop above] (w3) to node {$a, b$} (w3) ;

      \node[world](v1)[right of=w1, xshift=1.25in, label={right:$p, \lnot q$}]{$w'_1$};
      \node[world](v3)[above of=v1,label={right:$p,q$}]{$w'_3$};
      \draw[<->] (v1) to node {$a$} (v3);
      \draw[->,loop below] (v1) to node {$a,b$} (v1);
      \draw[->,loop above] (v3) to node {$a,b$} (v3) ;

      \node(m)[below=of w1]{$\mModel M$};
      \node(m')[below=of v1]{$\mModel M \mid p$};

      \draw[-,dotted] (w1) to node {\footnotesize 宣告 $p$}(v1) ;

    \end{tikzpicture}
  \end{center}
  \caption{宣告 $p$ 之前与之后。}
  \ollabel{fig:announcement-example}
\end{figure}

我们可以将公式的公共宣告视为使模型收缩，或者说将其限制到该公式为真的那些世界。\olref{fig:announcement-example} 给出了公开宣告~$p$ 的效果示例。该模型中值得注意的一点是，主体~$b$ 通过宣告得知了~$p$，而主体~$a$ 则没有（因为 $a$ 已经知道~$p$ 为真）。

更形式化地说，我们有 $\mSat{M}{\lnot \Knows_b p}[w_1]$ 但 $\mSat{M
\mid p}{\Knows_b p}[w'_1]$。这意味着 $\mSat{M}{[p] \Knows_b
p}[w_1]$。但关于宣告的结果，我们甚至可以提出一些更强的断言。事实上，$\mSat{M}{[p]\CKnows_{\{a,b\}} p}[w_1]$ 成立。换句话说，在 $p$ 被宣告之后，它成为了\emph{公共知识}。

我们或许会想，这一性质在一般情况下是否成立，以及如实地宣告~$!A$ 是否\emph{总是}会导致 $!A$ 成为公共知识。答案是否定的，这可能令人惊讶。事实上，确实可能存在一些公式，在如实地宣告它们之后，它们就不再为真了。例如，考虑在 \olref{fig:announcement-example} 的~$w_1$ 处宣告 $p \land \lnot \Knows_b p$ 的效果。事实上，$\mModel M \mid p$ 和 $\mModel M \mid (p \land \lnot \Knows_b p)$ 是同一个模型。
然而，正如我们已经注意到的，$\mSat{M \mid p}{\Knows_b
p}[w'_1]$。因此，$\mSat{M \mid (p \land \lnot \Knows_b p)}{\lnot
(p \land \lnot \Knows_b p)}[w'_1]$，所以这是一个一旦被宣告就会变为假的公式。
\end{document}